%% ============================================================
%% ENTREGABLE — TE3001B Fundamentación de Robótica
%% Teleoperación con Sensado de Fuerza: xArm 6 (Maestro-Esclavo)
%% Formato IEEE Transactions — Plantilla para Alumnos
%% Prof. Alberto Muñoz — Computational Robotics Lab
%% Tecnológico de Monterrey — 2026
%% ============================================================
\documentclass[journal,10pt,twoside]{IEEEtran}

%% -- Paquetes base --------------------------------------------
\UseRawInputEncoding
\usepackage[utf8]{inputenc}
\usepackage[T1]{fontenc}
\usepackage{placeins}
\usepackage[spanish,es-nodecimaldot]{babel}
\usepackage{amsmath,amssymb,bm}
\usepackage{graphicx} 
\usepackage{xcolor}
\usepackage{hyperref}
\usepackage{booktabs}
\usepackage{multirow}
\usepackage{array}
\usepackage{colortbl}
\usepackage{caption}
\usepackage{subcaption}
\usepackage{listings}
\usepackage{tcolorbox}
\usepackage{enumitem}
\usepackage{tikz}
\usepackage{pgfplots}
\usepackage{url}
\pgfplotsset{compat=1.18}
\usetikzlibrary{arrows.meta,shapes,positioning,calc,backgrounds}
\tcbuselibrary{breakable,skins}

\hypersetup{colorlinks=true, linkcolor=tec_blue, urlcolor=tec_blue,
            citecolor=tec_blue}

%% -- Colores institucionales ----------------------------------
\definecolor{tec_blue}{RGB}{0,68,124}
\definecolor{tec_lightblue}{RGB}{0,140,200}
\definecolor{tec_red}{RGB}{190,30,45}
\definecolor{tec_gray}{RGB}{100,100,100}
\definecolor{code_bg}{RGB}{245,247,250}
\definecolor{code_kw}{RGB}{0,100,180}
\definecolor{code_str}{RGB}{180,50,20}
\definecolor{code_cmt}{RGB}{80,140,80}
\definecolor{box_todo}{RGB}{255,243,205}
\definecolor{box_todo_border}{RGB}{200,140,0}
\definecolor{box_rubric}{RGB}{232,245,233}
\definecolor{box_rubric_border}{RGB}{67,160,71}
\definecolor{box_theory}{RGB}{227,242,253}
\definecolor{box_theory_border}{RGB}{25,118,210}

%% -- Cajas de tcolorbox --------------------------------------─
\newtcolorbox{todobox}[1][Completar]{
  colback=box_todo, colframe=box_todo_border,
  title={\textbf{$\triangleright$ #1}},
  fonttitle=\bfseries\color{tec_gray},
  breakable, enhanced, left=4pt, right=4pt
}
\newtcolorbox{rubricbox}[1][Rúbrica]{
  colback=box_rubric, colframe=box_rubric_border,
  title={\textbf{\checkmark\ Criterio de Evaluación — #1}},
  fonttitle=\bfseries\small\color{box_rubric_border},
  breakable, enhanced, left=4pt, right=4pt
}
\newtcolorbox{theorybox}[1][Marco Teórico]{
  colback=box_theory, colframe=box_theory_border,
  title={\textbf{(i) Contexto: #1}},
  fonttitle=\bfseries\small\color{tec_blue},
  breakable, enhanced, left=4pt, right=4pt
}

%% -- Estilo de código ----------------------------------------─
\lstdefinestyle{python}{
  language=Python,
  backgroundcolor=\color{code_bg},
  basicstyle=\ttfamily\scriptsize,
  keywordstyle=\color{code_kw}\bfseries,
  stringstyle=\color{code_str},
  commentstyle=\color{code_cmt}\itshape,
  numbers=left, numberstyle=\tiny\color{tec_gray},
  numbersep=6pt, frame=single, tabsize=4,
  showstringspaces=false, breaklines=true,
  morekeywords={import,from,as,def,class,return,True,False,None,self},
  literate={á}{{\'{a}}}1 {é}{{\'{e}}}1 {í}{{\'{i}}}1
           {ó}{{\'{o}}}1 {ú}{{\'{u}}}1 {ñ}{{\~{n}}}1,
}
\lstset{style=python}

%% -- Macros de ayuda ------------------------------------------
\newcommand{\answerline}[1][3cm]{\underline{\hspace{#1}}}
\newcommand{\pts}[1]{\textbf{\color{tec_red}[#1 pts]}}
\newcommand{\TODO}[1]{\begin{todobox}[Respuesta del equipo]
  \textit{#1}
\end{todobox}}
\newcommand{\FIGURE}[1]{\begin{todobox}[Insertar figura]
  \centering
  \begin{tikzpicture}
    \draw[tec_gray, dashed, line width=1pt] (0,0) rectangle (8,4);
    \node[tec_gray, font=\small] at (4,2) {#1};
  \end{tikzpicture}
\end{todobox}}

%% -- Metadatos del artículo ----------------------------------─
\begin{document}

\title{Teleoperación con Sensado de Fuerza Física\\
  entre dos xArm 6:\\
  Diseño, Implementación y Evaluación Experimental}

\author{%
  \IEEEauthorblockN{%
    Dante Alessandro Cortés Hernández  \IEEEauthorrefmark{1},\;
    Ana Karen Moreno Acevedo \IEEEauthorrefmark{1},\;
    Carlos Adriel Perez Silva	 \IEEEauthorrefmark{1},\;
    Ana Cecilia Vázquez Yañez 	\IEEEauthorrefmark{1}%
  }\\
  \IEEEauthorblockA{%
    \IEEEauthorrefmark{1}%
    Computational Robotics Lab, Tecnológico de Monterrey,
    Monterrey, N.L., México\\
    \texttt{\{A01255460, A00839936, A01773018, A00839986\}@tec.mx, }
  }
}

\maketitle

%% -- Abstract ------------------------------------------------─
\begin{abstract}
Este trabajo presenta el diseño e implementación de una arquitectura de teleoperación maestro--esclavo entre dos 
manipuladores xArm 6 con retroalimentación basada en contacto físico. En la propuesta, el robot esclavo ejecuta 
una trayectoria de aproximación e interacción con un dispositivo mecánico instrumentado con un sensor de fuerza, 
de modo que, al producirse contacto, la fuerza generada se mide y se transmite en tiempo real al sistema maestro 
mediante una red local. A partir de esta información, el robot maestro recibe una señal asociada al contacto 
ocurrido en el entorno remoto y modifica su comportamiento para representar una consecuencia análoga, como 
detención o retirada. De esta manera, el sistema no solo sincroniza el movimiento entre ambos brazos, sino que 
también comparte información de interacción con el entorno, planteando una forma práctica de teleoperación bidireccional
basada en comportamiento. La arquitectura integra sensado de fuerza en el esclavo, comunicación entre computadoras de 
control y una estrategia de respuesta cinemática en el maestro. El objetivo del trabajo es construir una plataforma 
funcional y reproducible para estudiar principios de contacto, transmisión de fuerza y retroalimentación entre 
manipuladores robóticos.
\end{abstract}

\begin{IEEEkeywords}
teleoperación, sensado de fuerza, xArm 6, maestro--esclavo, contacto físico, interacción con el entorno
\end{IEEEkeywords}

%% ============================================================
\section{Introducción}
%% ============================================================

\IEEEPARstart{E}{n} los sistemas de teleoperación, la utilidad del enlace maestro--esclavo aumenta significativamente
cuando la interacción remota no se limita al seguimiento de trayectorias, sino que también incorpora información 
del contacto con el entorno. En este contexto, la retroalimentación de fuerza permite que un evento mecánico 
ocurrido en el sistema remoto tenga una representación significativa en el sistema local. Esto resulta 
especialmente valioso en tareas donde la detección de contacto, la compresión de objetos o la interacción 
con superficies no puede tratarse únicamente como un fenómeno geométrico.

El xArm 6 ofrece una plataforma adecuada para estudiar este tipo de problemas, ya que su arquitectura de seis grados 
de libertad y su interfaz de programación permiten implementar un esquema maestro--esclavo en el que el movimiento 
y la información de interacción pueden intercambiarse entre dos estaciones de trabajo conectadas en red.

\begin{theorybox}[Motivación]
La idea central de este trabajo es que el contacto ocurrido en el robot esclavo no permanezca aislado en el entorno 
remoto. Cuando el esclavo comprime un dispositivo físico basado en resorte, la fuerza asociada a esa interacción 
se mide y se transmite al robot maestro para que éste represente una consecuencia equivalente del contacto. 

En lugar de implementar una reflexión háptica ideal basada en torque directo, el proyecto adopta una estrategia 
práctica donde la información de fuerza modifica el comportamiento del maestro mediante una respuesta observable, 
como detención, reducción de avance o retirada.
\end{theorybox}

\subsection{Contexto y Motivación}

La teleoperación con retroalimentación de fuerza es fundamental en la robótica moderna para lograr la transparencia 
en tareas de interacción física remota. Su relevancia radica en que permite al operador percibir no solo el movimiento,
sino también las fuerzas de contacto y la rigidez del entorno, lo que incrementa la precisión y la seguridad de la 
tarea. Esta tecnología es crítica en aplicaciones como la cirugía robótica mínimamente invasiva, donde el médico 
debe distinguir la resistencia de diferentes tejidos biológicos~\cite{sanchez2002}, y en la inspección de entornos 
hostiles, como centrales nucleares o mantenimiento submarino, donde el sensado de fuerza ayuda a evitar daños 
estructurales al interactuar con objetos frágiles o pesados~\cite{postigo2001}.

En el estado del arte se han propuesto distintos esquemas de control para gestionar esta interacción, 
destacando el control de impedancia, que modela la relación dinámica entre posición y fuerza y contribuye a mantener 
la estabilidad del sistema aun en presencia de retardos de red~\cite{postigo2001}. De forma complementaria, 
la incorporación de interfaces hápticas en sistemas con múltiples grados de libertad ha permitido una manipulación 
más intuitiva y una interacción más rica entre operador y entorno remoto~\cite{sanchez2002}. En todos estos 
enfoques, la idea central es que el sistema maestro no permanezca ajeno al contacto ocurrido en el esclavo, 
sino que reciba información suficiente para representar sus efectos.

Bajo esta perspectiva, en este proyecto se plantea una implementación práctica con dos xArm 6, donde el robot 
esclavo interactúa con un sensor mecánico basado en un resorte. La fuerza generada durante dicho contacto se 
captura y se transmite al sistema maestro mediante una red local entre computadoras. En lugar de depender 
necesariamente de una reproducción ideal de torque reflejado, la retroalimentación se plantea como una 
respuesta cinemática de retirada o modificación de movimiento en el robot maestro, con el fin de imitar 
la consecuencia mecánica del contacto detectado por el esclavo. Este enfoque permite estudiar de manera 
accesible los principios de teleoperación bidireccional con realimentación de fuerza~\cite{cardenas2016}.

\subsection{Contribuciones del Trabajo}

Las contribuciones principales de este trabajo son las siguientes:

\begin{itemize}
    \item Se presenta una arquitectura maestro--esclavo entre dos xArm 6 capaz de intercambiar información de movimiento
    y de interacción con el entorno en tiempo real.
    
    \item Se define una tarea de contacto en la que el robot esclavo interactúa con un sensor físico basado en un 
    resorte, con el propósito de capturar la fuerza aplicada durante la compresión.
    
    \item Se implemneta un mecanismo de transmisión de datos desde el esclavo hacia el maestro, de 
    modo que la información de fuerza obtenida en el contacto pueda ser enviada al segundo robot a través de la red.
    
    \item Se propone una estrategia de retroalimentación basada en comportamiento, en la que el robot 
    maestro modifica su movimiento para representar una consecuencia análoga al contacto detectado por el esclavo.
    
    \end{itemize}
      
%% ============================================================
\section{Marco Teórico}
\label{sec:background}
%% ============================================================

\subsection{Cinemática y Jacobiano del xArm 6}

El xArm 6 es un manipulador serial de $n = 6$ grados de libertad con articulaciones rotacionales. Su cinemática directa
puede describirse mediante la convención de Denavit--Hartenberg (DH) o el formalismo de Product of Exponentials (PoE).

\begin{theorybox}[Jacobiano Geométrico]
El Jacobiano $\bm{J}(\bm{q}) \in \mathbb{R}^{6\times6}$ relaciona las velocidades articulares con la velocidad 
cartesiana del efector final:
\[
\begin{pmatrix} \bm{v} \\ \bm{\omega} \end{pmatrix} = \bm{J}(\bm{q})\,\dot{\bm{q}}
\]
donde $\bm{v}$ es la velocidad lineal y $\bm{\omega}$ la velocidad angular. Su transpuesta mapea fuerzas 
cartesianas a torques articulares:
\[
\bm{\tau} = \bm{J}^\top(\bm{q})\,\bm{F}_{ext}
\]
Esta relación es la base del sensado de fuerza por torque articular.
\end{theorybox}

En este proyecto la fuerza no se estimó por torque articular sino mediante una celda de carga HX711 
montada en el efector del esclavo. La formulación del 
Jacobiano se presenta como fundamento teórico del análisis fuerza--movimiento.

\subsubsection*{Pregunta~T1}

\begin{enumerate}[label=(\alph*)]

\item El Jacobiano geométrico relaciona velocidades articulares con la velocidad lineal 
y angular física del efector:
\[
\begin{pmatrix} \bm{v} \\ \bm{\omega} \end{pmatrix} = \bm{J}_g(\bm{q})\dot{\bm{q}}
\]
El Jacobiano analítico relaciona velocidades articulares con la derivada de una parametrización 
de la pose, como ángulos de Euler. La diferencia es que el analítico depende de la representación 
de orientación elegida, mientras que el geométrico trabaja con velocidad angular física. Se usó el 
Jacobiano geométrico porque su transpuesta conecta fuerzas cartesianas con torques articulares de forma directa:
\[
\bm{\tau} = \bm{J}^\top(\bm{q})\bm{F}_{ext}
\]

\item Dos singularidades del xArm 6:
\begin{itemize}
    \item Singularidad de codo estirado: brazo y antebrazo quedan alineados. El Jacobiano pierde rango y el robot 
    pierde movilidad en ciertas direcciones cartesianas.
    \item Singularidad de muñeca: dos ejes de rotación de la muñeca se alinean. El determinante del Jacobiano 
    tiende a cero y la matriz se vuelve mal condicionada.
\end{itemize}
En ambos casos la estimación $\hat{\bm{F}} = (\bm{J}^\top)^{\dagger}\bm{\tau}$ se vuelve poco confiable porque 
la pseudoinversa amplifica el ruido en los torques medidos.

\begin{figure}[h]
    \centering
    \includegraphics[width=0.78\columnwidth]{t1_codo}
    \caption{Singularidad de codo estirado.}
    \label{fig:t1_codo}
\end{figure}

\begin{figure}[h]
    \centering
    \includegraphics[width=0.78\columnwidth]{t1_muneca}
    \caption{Singularidad de muñeca.}
    \label{fig:t1_muneca}
\end{figure}

\item La cercanía a una singularidad se detecta con el número de condición del Jacobiano:
\[
\kappa(\bm{J}) = \frac{\sigma_{\max}(\bm{J})}{\sigma_{\min}(\bm{J})}
\]
El robot está cerca de una singularidad cuando $\kappa(\bm{J}) > \kappa_{th}$, o equivalentemente 
cuando $\sigma_{\min}(\bm{J}) < \varepsilon$.

\end{enumerate}

\begin{figure}[h]
    \centering
    \includegraphics[width=0.92\columnwidth]{t1_flujo}
    \caption{Flujo de estimación de fuerza y control de singularidad.}
    \label{fig:t1_flujo}
\end{figure}

\FloatBarrier

\subsection{Dinámica del Manipulador}

La dinámica de un robot rígido de $n$ grados de libertad puede
expresarse como

\begin{equation}
\bm{M}(\bm{q})\ddot{\bm{q}} + \bm{C}(\bm{q},\dot{\bm{q}})\dot{\bm{q}}
+ \bm{g}(\bm{q}) = \bm{\tau} - \bm{J}^\top\bm{F}_{ext}
\label{eq:dynamics}
\end{equation}

donde $\bm{M}$ es la matriz de inercia del manipulador,
$\bm{C}$ representa los términos de Coriolis y centrífugos
y $\bm{g}$ corresponde al vector de gravedad.

\subsubsection*{Resolución de la Pregunta T2}

\begin{enumerate}[label=(\alph*)]

\item Considerando el caso nominal sin perturbación externa
($\bm{F}_{ext} \approx \bm{0}$), la ley de control de par calculado
(Computed Torque Control, CTC) se define como

\[
\bm{\tau} =
\bm{M}(\bm{q})\bm{a}_d
+
\bm{C}(\bm{q},\dot{\bm{q}})\dot{\bm{q}}
+
\bm{g}(\bm{q})
\]

donde la aceleración auxiliar está dada por

\[
\bm{a}_d =
\ddot{\bm{q}}_d
+
\bm{K}_v(\dot{\bm{q}}_d - \dot{\bm{q}})
+
\bm{K}_p(\bm{q}_d - \bm{q})
\]

Sustituyendo la ley de control en la Ec.~\eqref{eq:dynamics} y
cancelando términos comunes se obtiene

\[
\ddot{\bm{q}} = \bm{a}_d
\]

Definiendo el error de seguimiento

\[
\bm{e} = \bm{q}_d - \bm{q},
\qquad
\dot{\bm{e}} = \dot{\bm{q}}_d - \dot{\bm{q}},
\]

y reordenando en términos del error se obtiene

\[
\ddot{\bm{e}} + \bm{K}_v\dot{\bm{e}} + \bm{K}_p\bm{e} = \bm{0}.
\]

Esta ecuación muestra que el error de seguimiento queda gobernado
por una dinámica lineal de segundo orden desacoplada por
articulación. La estabilidad asintótica global del sistema se
garantiza si las matrices de ganancia $\bm{K}_p$ y $\bm{K}_v$
son simétricas definidas positivas, lo que implica rigidez
y amortiguamiento positivos en cada articulación.

\item El bucle de control interno del xArm 6 en modo
\texttt{SERVO} de la SDK opera aproximadamente a
$\mathbf{1.5\,kHz}$. En nuestra implementación el
controlador externo trabaja a una frecuencia de
$100\,\text{Hz}$ ($DT = 0.01\,\text{s}$).

Esta diferencia de frecuencias impone una limitación
práctica en la selección de las ganancias $\bm{K}_p$
y $\bm{K}_v$. Si las ganancias se incrementan demasiado,
el sistema puede volverse sensible al muestreo,
al retardo computacional y al ruido de medición,
lo que puede provocar oscilaciones o saturación del
esfuerzo de control. Por esta razón, el ancho de banda
del lazo cerrado debe mantenerse por debajo de la
frecuencia de actualización del controlador.

\item Con las ganancias utilizadas en el controlador

\[
\bm{K}_p = \mathrm{diag}(120,\,100,\,80),
\qquad
\bm{K}_v = \mathrm{diag}(25,\,20,\,15)
\]

la ecuación característica para cada articulación
puede escribirse como

\[
s^2 + k_{v,i}s + k_{p,i} = 0.
\]

Comparando con la forma estándar de un sistema de
segundo orden

\[
s^2 + 2\zeta_i\omega_{n,i}s + \omega_{n,i}^2 = 0,
\]

se obtiene

\[
\omega_{n,i} = \sqrt{k_{p,i}},
\qquad
\zeta_i = \frac{k_{v,i}}{2\sqrt{k_{p,i}}}.
\]

De esta forma,

\[
\zeta_1 \approx 1.14,
\qquad
\zeta_2 = 1.00,
\qquad
\zeta_3 \approx 0.84.
\]

Estos valores indican una sintonización cercana al
amortiguamiento crítico. En consecuencia,
puede considerarse un valor representativo de

\[
\zeta \approx 1,
\]

lo que proporciona una respuesta rápida,
con sobreimpulso mínimo y un seguimiento
estable de la trayectoria.

\end{enumerate}

\subsection{Control de Impedancia}

El control de impedancia modela el efector final como un sistema
masa--resorte--amortiguador virtual, de modo que la interacción
entre posición y fuerza pueda regularse de manera dinámica. Su
forma general puede escribirse como

\begin{equation}
\bm{F}_{ctrl}
=
\bm{K}_d(\bm{x}_d-\bm{x})
+
\bm{B}_d(\dot{\bm{x}}_d-\dot{\bm{x}})
+
\bm{M}_d(\ddot{\bm{x}}_d-\ddot{\bm{x}})
\label{eq:impedance}
\end{equation}

donde $\bm{K}_d$ representa la rigidez virtual, $\bm{B}_d$ el
amortiguamiento virtual y $\bm{M}_d$ la masa virtual del sistema.

\vspace{5mm}

\begin{enumerate}[label=(\alph*)]

\item El aumento de la rigidez virtual $\bm{K}_d$ tiene tres efectos
principales en el sistema de teleoperación. En primer lugar,
incrementa la sensación de rigidez percibida por el operador, ya
que pequeñas desviaciones entre la posición deseada y la posición
real generan fuerzas de interacción mayores. En segundo lugar,
una rigidez elevada reduce el margen de estabilidad cuando existe
retardo de red $T_d$, porque el sistema se vuelve más sensible a
errores temporales y a oscilaciones inducidas por la comunicación.
En tercer lugar, un valor grande de $\bm{K}_d$ acelera la respuesta
ante contacto, ya que el efector reacciona con mayor intensidad
frente a una penetración o desplazamiento en el entorno.

Desde el punto de vista práctico, esto implica un compromiso: una
rigidez baja produce una sensación más suave pero menos precisa,
mientras que una rigidez alta mejora la respuesta local al contacto
pero incrementa el riesgo de oscilación o comportamiento no pasivo.
Por esta razón, en teleoperación con red y manipuladores físicos no
conviene maximizar $\bm{K}_d$ sin considerar el retardo total del
sistema.

\item Para analizar el límite de rigidez con retardo de red se empleó
el criterio de pasividad de Colgate,

\[
K_{d,\max}
<
\frac{\pi}{2\,T_d\,\omega_c}.
\]

En el laboratorio se consideró un retardo de red representativo de

\[
T_d = 20~\text{ms} = 0.02~\text{s}.
\]

Para la selección de $\omega_c$ se adoptó una aproximación
conservadora basada en el lazo efectivo de teleoperación completo,
no solamente en la electrónica interna del xArm. Aunque la SDK del
robot trabaja con un servo interno rápido, la dinámica global del
sistema queda limitada por la actualización externa, la red, el
procesamiento del sensor de fuerza y la ejecución de comandos
cartesianos. Por ello, se tomó una banda efectiva de

\[
f_c = 10~\text{Hz},
\qquad
\omega_c = 2\pi f_c = 2\pi(10) \approx 62.83~\text{rad/s}.
\]

Sustituyendo en el criterio anterior,

\[
K_{d,\max}
<
\frac{\pi}{2(0.02)(62.83)}
\approx 1.25.
\]

Este resultado indica que, bajo una interpretación estricta de
pasividad con $T_d=20$~ms, la rigidez virtual admisible para una
retroalimentación háptica ideal es muy reducida. En consecuencia,
una implementación agresiva de impedancia o de reflexión de fuerza
directa habría comprometido la estabilidad del sistema.

Por esta razón, en el presente proyecto no se implementó una
impedancia háptica rígida en el maestro. En su lugar, se adoptó una
estrategia práctica de retroalimentación basada en comportamiento,
en la que la fuerza medida por la celda de carga en el esclavo
provoca en el maestro una consecuencia observable, como detención,
reducción de avance o retirada, manteniendo así una interacción
estable y segura.

\item Desde el punto de vista conceptual, el esquema de impedancia
empleado para interpretar la interacción se formuló en
\textbf{espacio cartesiano}, ya que la tarea de teleoperación se
define en términos de movimiento del efector final y de contacto
con el entorno. En este contexto, resulta más natural expresar la
relación fuerza--desplazamiento en las coordenadas cartesianas
$\bm{x}$ del efector que en las coordenadas articulares
$\bm{q}$ del robot.

Además, la variable física de interés en el experimento es la fuerza
de contacto medida en el extremo del manipulador esclavo mediante una
celda de carga con HX711, por lo que su interpretación es
directamente cartesiana. Modelar la interacción en espacio
articular habría complicado la relación entre la fuerza medida, la
deformación del resorte y la respuesta observable del sistema.
Por esta razón, la elección de un marco cartesiano resulta la más
adecuada para la tarea de teleoperación desarrollada.

\end{enumerate}
\vspace{5mm}

Con base en lo anterior, puede concluirse que el control de
impedancia constituye un marco teórico útil para analizar la
interacción fuerza--movimiento del sistema. Sin embargo, debido al
retardo de red y a las limitaciones prácticas de la plataforma, la
implementación experimental se orientó a una respuesta cinemática
estable basada en la fuerza medida, en lugar de una reflexión
háptica ideal de alta rigidez.

\begin{table}[!t]
\centering
\caption{Efecto cualitativo de distintos valores de rigidez virtual.}
\label{tab:kd_compare}
\small
\begin{tabular}{cccc}
\toprule
\textbf{$K_d$} & \textbf{Sensación} & \textbf{Respuesta} & \textbf{Estabilidad} \\
\midrule
Bajo   & Suave   & Lenta  & Alta \\
Medio  & Balanceada & Moderada & Buena \\
Alto   & Rígida  & Rápida & Reducida \\
\bottomrule
\end{tabular}
\end{table}

%% ============================================================
\section{Arquitectura del Sistema}
\label{sec:architecture}
%% ============================================================

\subsection{Visión General}
El sistema consiste en dos manipuladores xArm6, cada uno conectado a su
propia PC de control dentro de una red local. El maestro envía referencias
de movimiento al esclavo, y el esclavo devuelve información de fuerza al
maestro. Ambos flujos ocurren en tiempo real.
El esclavo ejecuta una trayectoria de aproximación hacia una celda de carga.
Al hacer contacto, la fuerza de reacción se mide y se transmite
al sistema maestro. Aunque el maestro no está en contacto físico con el
entorno, recibe esa fuerza y reacciona a ella modificando su trayectoria
y estado dinámico.
La arquitectura busca compartir el efecto del contacto entre ambos brazos.
El sistema incorpora una capa de retroalimentación proveniente del entorno,
lo que le da a la teleoperación un carácter bidireccional. La compresión
del resorte en el esclavo se refleja directamente en el comportamiento
del maestro.

\subsubsection*{Pregunta~A1}

\begin{enumerate}[label=(\alph*)]

\item La arquitectura se organiza como un sistema maestro--esclavo distribuido en dos PCs conectadas por red local.
La PC Maestra coordina el xArm 6 Maestro y ejecuta un modelo cinemático-dinámico planar de 3 GDL para calcular la 
referencia cartesiana e incorporar la retroalimentación de fuerza del esclavo. La PC Esclava recibe esa referencia, 
la convierte en comandos cartesianos para el xArm 6 Esclavo y procesa la señal del sensor HX711 recibida por UDP.

Las señales que fluyen entre bloques son:
\begin{itemize}
    \item Maestro a Esclavo: posición cartesiana deseada $\bm{x}_d \in \mathbb{R}^2$ y estado de pinza $b \in \{0,1\}$.
    \item Esclavo a Maestro: fuerza de contacto $\bm{F}_e \in \mathbb{R}^2$ y bandera de contacto $c \in \{0,1\}$.
    \item ESP32 a PC Esclava: fuerza escalar $F \in \mathbb{R}$ medida por el HX711.
\end{itemize}

\begin{figure}[h]
    \centering
    \includegraphics[width=0.9\columnwidth]{diagrama_bloques}
    \caption{Diagrama de bloques del sistema de teleoperación maestro-esclavo.}
    \label{fig:diagrama_bloques}
\end{figure}

\item En cada ciclo de control el maestro calcula $\bm{x}_m \in \mathbb{R}^2$ a partir del estado 
articular $\bm{q}_m \in \mathbb{R}^3$ y lo envía al esclavo por UDP. En paralelo, publica velocidades cartesianas 
al xArm físico via \texttt{/servo\_server/delta\_twist\_cmds}. El esclavo recibe $\bm{x}_d$, lo convierte en un 
comando para el xArm real con \texttt{set\_servo\_cartesian\_aa}, lee la fuerza del HX711 y la devuelve al 
maestro como $\bm{F}_e$. El maestro incorpora $\bm{F}_e$ en su ley de control, cerrando el lazo bilateral.

\item La frecuencia de actualización es $f_s = 100\,\text{Hz}$ ($DT = 0.01\,\text{s}$). 
Se eligió \texttt{set\_servo\_cartesian\_aa} en el esclavo y MoveIt Servo con mensajes \texttt{TwistStamped} 
en el maestro porque ambos aceptan comandos continuos a alta frecuencia en espacio cartesiano. 
\texttt{set\_position} no es adecuado para teleoperación porque introduce latencia y está orientado a 
movimientos discretos.

\end{enumerate}

\subsubsection*{Pregunta~A2}

\begin{enumerate}[label=(\alph*)]

\item Se eligió UDP. No garantiza entrega ni orden,
pero elimina la sobrecarga de TCP. En control a 100~Hz, un paquete retrasado
es peor que uno perdido.

\begin{center}
\small
\begin{tabular}{lcc}
  \toprule
  \textbf{Métrica} & \textbf{UDP} & \textbf{TCP} \\
  \midrule
  RTT promedio (ms)      & 1.92  & $\sim$4--8 \\
  Jitter (ms)            & 0.78  & $\sim$2--5 \\
  Paquetes perdidos (\%) & 0.0   & N/A \\
  Throughput (KB/s)      & $>$10 & $>$10 \\
  \bottomrule
\end{tabular}
\end{center}

\item Dos tipos de paquetes circulan en el sistema:

Maestro a Esclavo: \texttt{\{"xd": [float, float], "gripper": int\}}, aprox. 40 bytes.

ESP32 a Esclavo: \texttt{\{"F": float\}}, aprox. 15 bytes al puerto 9003.

Sin header ni checksum. La red LAN mostró 0\% de pérdida, por lo que no se necesita.

\item Con $T_{RTT} = 1.92$~ms:
\[
  T_d = \frac{1.92}{2} = 0.96 \text{ ms}
\]
Aplicando el criterio de Colgate con $\omega_c = 2\pi \cdot 10$~rad/s:
\[
  K_{d,max} < \frac{\pi}{2 \cdot 0.00096 \cdot 62.83} \approx 25{,}900 \text{ N/m}
\]
El sistema opera con $K_d = 500$~N/m, muy por debajo del límite. El retardo
medido no compromete la estabilidad háptica.

\end{enumerate}

\begin{figure}[h]
    \centering
    \includegraphics[width=0.8\columnwidth]{nettest}
    \caption{Resultados de net\_test.py entre PC Maestro y PC Esclavo.}
    \label{fig:nettest}
\end{figure}

%% ============================================================
\section{Estrategia de Sensado de Fuerza}
\label{sec:force}
%% ============================================================

\begin{theorybox}[¿Cómo «siente» la fuerza un xArm sin sensor F/T externo?]
  El xArm 6 permite leer los torques articulares a través de
  \texttt{arm.get\_joint\_states()} o mediante el modo de control de
  torque. A partir de los torques medidos $\bm{\tau}_{meas}$ y el
  modelo dinámico del robot, es posible \emph{estimar} la fuerza
  externa en el efector final sin instalar un sensor F/T dedicado.
  Esta técnica se denomina \textbf{sensado de fuerza basado en torque articular}.
\end{theorybox}

\subsection{Estimación de Fuerza por Torque Articular}

La fuerza externa en el efector se estima como:
\begin{equation}
  \hat{\bm{F}}_{ext} = \left(\bm{J}^\top\right)^+
    \underbrace{\left(\bm{\tau}_{meas} - \bm{\tau}_{model}\right)}_{\Delta\bm{\tau}}
  \label{eq:force_estimate}
\end{equation}
donde $\bm{\tau}_{model} = \bm{M}(\bm{q})\ddot{\bm{q}} +
\bm{C}\dot{\bm{q}} + \bm{g}(\bm{q})$ es el torque predicho por el modelo.

\subsubsection*{Pregunta~F1}

\begin{enumerate}[label=(\alph*)]

\item Tres fuentes de error en la estimación de $\hat{\bm{F}}_{ext}$:
\begin{itemize}
    \item Error de modelo dinámico: el modelo $\bm{M}$, $\bm{C}$, $\bm{g}$
    es una aproximación. Parámetros inerciales incorrectos generan errores
    de hasta 1--2~N en configuraciones alejadas del home.
    \item Fricción articular: el modelo no incluye fricción en las
    articulaciones. En el xArm6, esto introduce un offset constante en
    $\Delta\bm{\tau}$ de aproximadamente 0.5--1~Nm por articulación.
    \item Ruido en los sensores de torque: los encoders y estimadores
    internos del xArm6 tienen ruido de cuantización. En reposo, esto
    se traduce en una fuerza estimada de aproximadamente 0.3--0.8~N.
\end{itemize}

\item La compensación de gravedad se realiza dentro del control de
impedancia. La expresión usada es:
\[
\bm{g}(\bm{q}) =
\begin{pmatrix}
\begin{aligned}
G\bigl(& (M_1+M_2+M_3)L_1c_1 \\
       & + (M_2+M_3)L_2c_{12} \\
       & + M_3L_3c_{123} \bigr)
\end{aligned}
\\[6pt]
\begin{aligned}
G\bigl(& (M_2+M_3)L_2c_{12} \\
       & + M_3L_3c_{123} \bigr)
\end{aligned}
\\[6pt]
G\,M_3L_3c_{123}
\end{pmatrix}
\]
donde $c_1=\cos(q_1)$, $c_{12}=\cos(q_1+q_2)$ y
$c_{123}=\cos(q_1+q_2+q_3)$. Los parámetros $M_i$ y $L_i$
se obtienen del modelo simplificado de masa puntual del robot.

\item Se aplica un filtro pasa-bajas de primer orden con frecuencia
de corte $f_c = 10$~Hz. La función de transferencia discreta es:
\[
H(z) = \frac{1-\alpha}{1 - \alpha z^{-1}}, \qquad
\alpha = e^{-2\pi f_c \, DT} \approx 0.535
\]
con $DT = 0.01$~s. Se eligió $f_c = 10$~Hz porque el movimiento
del robot opera por debajo de esa frecuencia, mientras que el ruido
de los sensores se concentra en frecuencias más altas.

\item Se usó un umbral $F_{th} = 2$~N. Con el robot en reposo, el
ruido de la fuerza estimada tiene una desviación estándar de
aproximadamente 0.3--0.5~N. El umbral de 2~N garantiza un SNR mayor
a 4, evitando falsas detecciones de contacto.

\end{enumerate}

\subsection{Alternativas de Sensado Consideradas}

\subsubsection*{Pregunta~F2}

\begin{enumerate}[label=(\alph*)]

\item Comparativa de estrategias de sensado:

\begin{center}
\footnotesize
\setlength{\tabcolsep}{3pt}
\renewcommand{\arraystretch}{1.0}
\begin{tabular}{p{2.2cm}p{2.1cm}p{2.2cm}p{1.3cm}}
  \toprule
  \textbf{Estrategia} & \textbf{Ventajas} & \textbf{Desventajas} & \textbf{Costo} \\
  \midrule
  Torque articular
    & Sin hardware adicional, usa sensores internos
    & Depende del modelo dinámico, baja precisión
    & Bajo \\
  Sensor F/T en muñeca
    & Medición directa y precisa
    & Alto costo, requiere calibración
    & Alto \\
  Visión por computadora
    & Sin contacto físico adicional
    & Lento, sensible a iluminación
    & Medio \\
  Galgas extensométricas
    & Alta sensibilidad, costo moderado
    & Instalación y calibración complejas
    & Medio \\
  \bottomrule
\end{tabular}
\end{center}

\item Se usaron galgas extensométricas (celda de carga HX711) en el
efector del robot esclavo. Esta elección se basó en tres criterios:
costo moderado, integración directa sin necesidad de sensores
comerciales de fuerza/torque, y medición local en el punto de contacto.

\item Con un sensor F/T de seis ejes en la muñeca, los cambios serían:
en hardware, montar el sensor entre la brida del robot y el efector
final. En software, leer la wrench del sensor via ROS2
(\texttt{geometry\_msgs/WrenchStamped}), transformar las fuerzas al
frame de la base con la cinemática directa, y reemplazar
\texttt{self.net.Fe} con los valores del sensor en lugar de los del
HX711. Esto permitiría medir fuerzas en los seis ejes en lugar de
solo la componente vertical.

\end{enumerate}

%% ============================================================
\section{Esquema de Control}
\label{sec:control}
%% ============================================================

\subsection{Control del Robot Esclavo}

\subsubsection*{Pregunta~C1}

\begin{enumerate}[label=(\alph*)]

\item La ley de control implementada en el esclavo es control de impedancia
en espacio cartesiano:
\[
\bm{F}_{ctrl} = \bm{K}_d(\bm{x}_d - \bm{x}) + \bm{B}_d(\dot{\bm{x}}_d - \dot{\bm{x}})
\]
\[
\bm{\tau} = \bm{J}^\top \bm{F}_{ctrl} + \bm{g}(\bm{q}) + \bm{C}(\bm{q},\dot{\bm{q}})\dot{\bm{q}}
\]
con $\bm{K}_d = 500$~N/m y $\bm{B}_d = 70$~N$\cdot$s/m.
Se usó \texttt{set\_servo\_cartesian\_aa} (modo 1 de la SDK). Este modo
acepta comandos cartesianos a alta frecuencia sin planificación de trayectoria,
lo cual es necesario para seguimiento en tiempo real a 100Hz.

\item El ciclo de control del esclavo opera como sigue:

\begin{lstlisting}[caption={Ciclo de control del robot Esclavo},
                   label={lst:slave}]
from xarm.wrapper import XArmAPI

arm = XArmAPI('192.168.1.195')
arm.motion_enable(enable=True)
arm.set_mode(1)   # modo servo cartesiano
arm.set_state(0)

while running:
    # 1. Recibir posicion del maestro via UDP
    x_des = net.x_des.copy()

    # 2. Escalar y enviar al xArm fisico
    px = x_des[0] * 700.0 + 300.0
    py = x_des[1] * 700.0
    arm.set_servo_cartesian_aa(
        [px, py, 300.0, 0, 0, 0], speed=100, mvacc=1000)

    # 3. Calcular fuerza de contacto (modelo impedancia)
    F_contact, _, in_contact = contact_model.compute_contact_force(x_ef)

    # 4. Leer fuerza del sensor HX711 via UDP
    Fe = net.Fe   # actualizado por hilo ESP32

    # 5. Enviar fuerza al maestro
    net.send_force(Fe, in_contact)

    time.sleep(0.01)
\end{lstlisting}

\item Se aplicó un factor de escala entre el espacio de simulación
y el espacio real del xArm6:
\[
\bm{x}_{slave}^{ref} = \alpha \cdot \bm{x}_{master} + \bm{x}_{offset}
\]
con
\[
\alpha = \frac{630\ \text{mm}}{0.9\ \text{m}} \approx 700\ \text{mm/m}
\]
La simulación opera en un rango de $\pm 0.9$~m, mientras que el xArm6
tiene un alcance de $\sim 700$~mm. El offset $\bm{x}_{offset} = [300, 0, 300]$~mm
centra el espacio de trabajo en la posición home del robot.

\item Los torques se saturan a $\pm 20$~N$\cdot$m en cada articulación:
\[
\bm{\tau} = \text{clip}(\bm{\tau},\ -20,\ +20)\ \text{N$\cdot$m}
\]
Esto protege al robot de comandos de fuerza excesivos durante contacto
brusco o singularidades cercanas.

\end{enumerate}

\subsubsection*{Pregunta~C2}

\begin{enumerate}[label=(\alph*)]

\item $\beta$ es el factor de escala háptica. Representa cuánta fuerza
siente el operador en el maestro por cada Newton medido en el esclavo.
Un $\beta > 1$ amplifica la sensación de fuerza y un $\beta < 1$ la reduce.

El valor máximo estable se calcula con la condición de pasividad de Colgate:
\[
\beta_{max} < \frac{1}{2 \cdot T_d \cdot f_c}
= \frac{1}{2 \cdot 0.00096 \cdot 62.83} \approx 8.3
\]
Se usó $\beta = 0.3$ para garantizar estabilidad con margen amplio
y evitar saturación de torques en el maestro.

\item El xArm6 no permite aplicar torques articulares directamente
desde la SDK en modo normal. La alternativa
implementada fue reflejar la fuerza del esclavo como perturbación en
la referencia cartesiana del maestro via el Jacobiano transpuesto,
dentro de la simulación Python:
\[
\bm{\tau}_{fb} = \bm{J}_m^\top \cdot \beta \cdot \hat{\bm{F}}_{ext}
\]

\item Fragmento del feedback háptico en el maestro:

\begin{lstlisting}[caption={Feedback háptico en master\_robot.py},
                   label={lst:haptic}]
# En computed_torque(), antes de retornar tau:
if F_ext is not None and np.linalg.norm(F_ext) > 0.01:
    J = jacobian_3r(q)
    tau += J.T @ (0.3 * F_ext)   # beta = 0.3

tau = np.clip(tau, -20.0, 20.0)
\end{lstlisting}

\end{enumerate}

%% ============================================================
\section{Implementación con xArm SDK}
\label{sec:implementation}
%% ============================================================

\subsubsection*{Pregunta~I1}

\begin{enumerate}[label=(\alph*)]

\item Configuración de ambos robots:
\begin{itemize}
    \item xArm6 Maestro: IP \texttt{192.168.1.226}, modo 1, estado 0.
    \item xArm6 Esclavo: IP \texttt{192.168.1.195}, modo 1, estado 0.
    \item Velocidad máxima TCP: 100~mm/s.
    \item Aceleración máxima TCP: 1000~mm/s$^2$.
    \item Torque máximo por articulación: 20~N$\cdot$m (saturación en software).
\end{itemize}

\item Ambos robots se llevaron a posición home antes de iniciar.
El espacio de trabajo se sincronizó con un factor $\alpha = 700$~mm/m
sobre las coordenadas de la simulación, con offset $[300, 0, 300]$~mm.

\item Dos limitaciones encontradas:
\begin{itemize}
    \item El modo servo necesita comandos continuos. Si el hilo se
    interrumpe más de 500~ms, el robot entra en error.
    Se resolvió con reconexión automática por excepciones.
    \item \texttt{set\_servo\_cartesian\_aa} rechaza poses fuera del
    workspace. Se resolvió saturando coordenadas antes de enviarlas.
\end{itemize}

\item El manejo de errores se muestra en el código a continuación.

\end{enumerate}

\begin{lstlisting}[caption={Configuración inicial de los xArm6},
                   label={lst:setup}]
from xarm.wrapper import XArmAPI

def setup_robot(ip):
    arm = XArmAPI(ip)
    arm.motion_enable(enable=True)
    arm.set_mode(1)
    arm.set_state(0)
    arm.set_tcp_maxacc(1000)
    arm.set_servo_angle(
        angle=[0, 0, 0, 0, 0, 0], speed=20, wait=True)
    return arm

arm_master = setup_robot('192.168.1.226')
arm_slave  = setup_robot('192.168.1.195')

try:
    arm_slave.set_servo_cartesian_aa(
        [px, py, 300.0, 0, 0, 0], speed=100, mvacc=1000)
except Exception as e:
    print(f'[xArm error] {e}')
    arm_slave.clean_error()
    arm_slave.set_state(0)
\end{lstlisting}

%% ============================================================
\section{Resultados Experimentales}
\label{sec:experiments}
%% ============================================================

\subsection{Configuración Experimental}

\subsubsection*{Pregunta~E1}

\begin{enumerate}[label=(\alph*)]

\item Especificaciones del hardware utilizado:

\begin{center}
\scriptsize
\setlength{\tabcolsep}{4pt}
\renewcommand{\arraystretch}{1.0}
\begin{tabular}{lll}
  \toprule
  \textbf{Componente} & \textbf{PC Maestro} & \textbf{PC Esclavo} \\
  \midrule
  CPU      & i7-11800H @ 2.30~GHz & i3-1115G4 @ 3.00~GHz \\
  RAM      & 7.6~GB               & 20.0~GB \\
  OS       & Ubuntu 22.04.5 LTS   & Ubuntu 22.04 LTS \\
  Python   & 3.10.12              & 3.10.12 \\
  xArm SDK & xArm SDK 1.17.3      & xArm-Python-SDK + xarm\_ros2 \\
  \bottomrule
\end{tabular}
\end{center}

\item El operador movió el robot maestro con WASD hacia la posición
del agujero (HOLE) en la simulación. El robot esclavo replicó el
movimiento e interactuó con una celda de carga HX711 montada en
su efector. El contacto fue deliberado: se presionó el efector
contra la celda para generar una fuerza medible y verificar
el feedback háptico en el maestro.

\item Se realizaron 3 ensayos. El criterio de éxito fue:
error cartesiano menor a 10mm durante el seguimiento y
detección de fuerza de contacto mayor al umbral $F_{th} = 2$~N
al momento de presionar la celda de carga.

\end{enumerate}

\subsection{Estimación y Retroalimentación de Fuerza}

\subsubsection*{Pregunta~E2}
\begin{enumerate}[label=(\alph*)]
\item Durante el evento de contacto controlado, la fuerza estimada
en el eje $y$ alcanzó un máximo de 6.4~N al presionar el efector
contra la celda de carga. Los ejes $x$ y $z$ se mantuvieron
por debajo del umbral de ruido (0.5~N) al tratarse de contacto
unidireccional. La fuerza superó el umbral $F_{th} = 2$~N a los
0.8~s de iniciado el ensayo.

\item Se aplicó una fuerza de referencia $F_{ref} = 5$~N con una
báscula en la dirección $y$ del efector. Los resultados se muestran
en la Fig.~\ref{fig:e3b}. Los valores obtenidos fueron:
\begin{itemize}
    \item Error de estimación: $\epsilon = |5.0 - 4.7| = 0.3$~N.
    \item Sensibilidad: 0.82~N por N$\cdot$m de torque articular.
    \item Ruido en reposo: 0.35~N (desviación estándar sin contacto).
\end{itemize}

\begin{figure}[h]
    \centering
    \includegraphics[width=0.92\columnwidth]{slave_plots_1}
    \caption{Fuerzas de contacto y torques registrados durante el ensayo.
    Se observa el evento de contacto y la respuesta del sistema de impedancia.}
    \label{fig:e3b}
\end{figure}

\item La fuerza háptica en el esclavo fue proporcional a la fuerza
estimada en el maesrto con factor $\beta = 0.3$:
\[
F_{haptic} = 0.3 \cdot \hat{F}_{slave}
\]
Al registrar $\hat{F}_{slave} = 6.4$~N, el maestro recibió
$F_{haptic} \approx 1.9$~N, consistente con lo observado en
las gráficas de torque del maestro.

\item Para la configuración articular al instante de primer contacto
$\bm{q} = [0.6,\ -0.5,\ 0.1]$~rad, el torque esperado es:
\[
\bm{\tau}_{esp} = \bm{J}^\top(\bm{q})\,\hat{\bm{F}}_{ext}
= \bm{J}^\top \begin{pmatrix} 0 \\ 2.0 \end{pmatrix}
\approx \begin{pmatrix} 1.8 \\ 1.2 \\ 0.4 \end{pmatrix}\ \text{N$\cdot$m}
\]
Estos valores son consistentes en orden de magnitud con los torques
registrados en las gráficas al momento del contacto.

\end{enumerate}


\subsection{Análisis de Estabilidad del Sistema Háptico}

\subsubsection*{Pregunta~E3}

\begin{enumerate}[label=(\alph*)]

\item El sistema fue estable en ambos casos. En simulación, el contacto
con la pared (obstáculo con $K_{obs} = 120$~N/m) generó fuerzas de
reacción que el control de impedancia absorbió sin oscilaciones.
En el objeto real, el efector presionó la celda de carga y el sistema
convergió a una fuerza estable sin rebotes.

\item El valor teórico calculado fue $\beta_{max} \approx 8.3$.
Experimentalmente, el sistema comenzó a mostrar oscilaciones leves
en la simulación al alcanzar $\beta \approx 7$, lo cual es consistente
con el valor teórico considerando que el modelo asume condiciones ideales.
Se operó con $\beta = 0.3$ para garantizar estabilidad con margen amplio.

\item No se observó acoplamiento de inestabilidad durante contacto
simultáneo. Esto se debe a que el feedback háptico usa $\beta = 0.3$,
muy por debajo del límite de pasividad, lo cual amortigua cualquier
acoplamiento entre los dos lazos de fuerza antes de que genere
oscilaciones sostenidas.

\end{enumerate}

%% ============================================================
\section{Discusión}
\label{sec:discussion}
%% ============================================================

\subsubsection*{Pregunta~D1}

\begin{enumerate}[label=(\alph*)]

\item El sistema, aunque funcional, carga con tres limitaciones que
no se pueden ignorar:

\begin{itemize}
    \item El HX711 es un sensor uniaxial. Solo ve lo que pasa en el
    eje vertical, y todo lo que ocurre en los lados simplemente
    no existe para el. La solucion obvia es un sensor F/T de seis ejes
    en la muneca del esclavo, que abriria los ojos del sistema
    a la riqueza completa del contacto.

    \item El modelo dinamico es una simplificacion que le cuesta caro
    al estimador de fuerza. Ignorar la friccion articular introduce
    errores de hasta 1--2~N, que en un sistema haptico se sienten.
    Identificar los parametros de friccion del xArm6
    articulacion por articulacion resolveria esto de raiz.

    \item El sistema vive en 2D. La simulacion opera en un plano
    $(x,y)$ mientras que el xArm6 tiene 6~DOF esperando ser usados.
    Extender el control de impedancia a 6D con velocidades de torsion
    despertaria todo el potencial del hardware disponible.
\end{itemize}

\item Frente al estado del arte, la comparacion es honesta.
Ozdamar et al.~\cite{ozdamar2022} lograron algo bastante mas ambicioso:
cuatro robots Franka Emika Panda con stiffness variable en tiempo real,
algo que este sistema ni intenta al trabajar con $K_d$ fijo.
En cambio, en latencia el sistema es competitivo. Un estudio reciente
de teleoperacion quirurgica ~\cite{surgical2023} reporto errores de 2.4~mm en VLAN y
7.8~mm en WLAN con latencias de hasta 500~ms. Aqui se alcanzo menos
de 10~mm con 0.96~ms de retardo, lo cual, para una red local, es
dificil de superar. El talon de Aquiles sigue siendo la fuerza planar
y la ausencia de medicion en seis ejes.

\item Si el sistema tuviera que sobrevivir con $T_d = 200$~ms,
la rigidez virtual maxima colapsaria a:
\[
K_{d,max} < \frac{\pi}{2 \cdot 0.2 \cdot 62.83} \approx 125\ \text{N/m}
\]
Operar con $K_d = 500$~N/m seria una receta para la inestabilidad.
La solucion esta en las wave variables~\cite{niemeyer2008}, una tecnica
que transforma las senales de fuerza y velocidad en variables de onda
que preservan la pasividad del sistema sin importar cuanto retardo
haya en la red. Con eso, incluso una conexion transcontinental
dejaria de ser un problema.

\item La aplicacion que mas le quedaria al sistema, con las mejoras
indicadas, es cirugia robotica remota. Un cirujano en una ciudad
moviendo el maestro, sintiendo en sus manos la resistencia del tejido
que el esclavo toca a cientos de kilometros de distancia. Con sensor
F/T de seis ejes, wave variables y control en 6D, esto deja de ser
ciencia ficcion para procedimientos de baja complejidad. El costo
de implementacion estaria en el orden de USD~50,000--100,000,
contando dos xArm6, sensores F/T comerciales y la red dedicada
de baja latencia que todo esto requiere.

\end{enumerate}

%% ============================================================
\section{Conclusiones}
\label{sec:conclusions}
%% ============================================================
Se logró implementar un sistema de teleoperación maestro-esclavo funcional entre dos brazos
robóticos xArm6 de UFACTORY, haciendo uso de un sensor HX711 para la retroalimentación de fuerza
puesto en el robot maestro. Se logró un error cartesiano menor a 10~mm durante la operación del robot,
y los picos de fuerza de contacto registrados durante la presión contra la celda de carga
fueron consistentes con el modelo de impedancia implementado.
Al realizar este reto, uno de los desafíos más importantes enfrentados fue la sincronización de la simulación 
de Python con el stack de ROS2 con MoveIt Servo y la SDK del xArm6. En particular, el modo servo del robot requiere 
comandos a alta frecuencia para mantener un movimiento fluido, cualquier interrupción en el hilo de control provocaba 
que el robot entrara en estado de error. Esto se resolvió aislando el envío de comandos al xArm en el hilo principal 
de simulación y añadiendo manejo de excepciones para evitar caídas del sistema.
El aprendizaje más relevante de este proyecto fue entender un poco más el Jacobiano y lo importante que es para 
el mapeo de este tipo de movimientos. Permite mapear velocidades articulares a velocidades cartesianas, además permite
traducir fuerzas en el espacio de tarea directamente a torques articulares. Sin duda, uno de los aprendizajes core 
de este proyecto, al ser importante para el control de impedancia del esclavo y el feedback del maestro, 
donde la fuerza se refleja como torque en las articulaciones del operador a través del mismo Jacobiano.

%% ============================================================
\section*{Agradecimientos}
%% ============================================================

Agradecemos al Computational Robotics Lab del Tecnológico de Monterrey por el acceso al equipo y al espacio de trabajo,
así como al Prof. Alberto Muñoz por la guía académica y el planteamiento del reto. 
También se agradece el apoyo del laboratorio y de nuestros compañeros de clase durante la preparación del 
sistema maestro--esclavo y del sensor de fuerza empleado en este proyecto.

%% ============================================================
%% BIBLIOGRAFÍA
%% ============================================================
\begin{thebibliography}{10}

\bibitem{xarm_github} UFACTORY, ``xArm-Python-SDK,'' GitHub repository, 2024. [Online]. Available: \url{https://github.com/xArm-developer/xArm-Python-SDK}

\bibitem{siciliano2009} B. Siciliano, L. Sciavicco, L. Villani, G. Oriolo, \textit{Robotics: Modelling, Planning and Control}, Springer, 2009.

\bibitem{spong2006} M.W. Spong, S. Hutchinson, M. Vidyasagar, \textit{Robot Modeling and Control}, Wiley, 2006.

\bibitem{hogan1985} N. Hogan, ``Impedance Control: An Approach to Manipulation,'' \textit{ASME J. Dyn. Sys., Meas., Control}, vol. 107, 1985.

\bibitem{colgate1993} J.E. Colgate, J.M. Brown, ``Factors affecting the Z-width of a haptic display,'' in \textit{Proc. IEEE ICRA}, 1994, pp. 3205--3210.

\bibitem{murray1994} R.M. Murray, Z. Li, S.S. Sastry, \textit{A Mathematical Introduction to Robotic Manipulation}, CRC Press, 1994.

\bibitem{niemeyer2008} G. Niemeyer, C. Preusche, G. Hirzinger, ``Telerobotics,'' in \textit{Springer Handbook of Robotics}, Springer, 2008.

\bibitem{thompson2001} S. Thompson, I. Reid, A. Mu\~noz, D.W. Murray, ``Using a stationary camera to improve the performance of a teleoperated robot,'' \textit{IEEE Trans. Systems, Man \& Cybernetics-A}, vol. 31, no. 1, pp. 74--82, 2001.

\bibitem{sanchez2002} E.J. Sánchez Tapia, ``Implementación de un sistema teleoperado con reflexión de fuerza de seis grados de libertad,'' Tesis Doctoral, Universidad de Navarra, 2002.

\bibitem{postigo2001} J.F. Postigo, V.A. Mut, R.O. Carelli, L.A. Baigorria, B.R. Kuchen, ``Teleoperación de Robots con Reflexión de Fuerza y Control de Impedancia,'' \textit{Computación y Sistemas}, vol. 4, no. 4, pp. 340--355, 2001.

\bibitem{cardenas2016} D.D. Cárdenas Salas, ``Teleoperación con realimentación de fuerza aplicada a un robot móvil con restricción no-holonómica,'' Proyecto de Titulación, Escuela Politécnica Nacional, Ecuador, 2016.

\bibitem{ref9} R.V. Patel, S.F. Atashzar, M. Tavakoli, ``Haptic Feedback and Force-Based Teleoperation in Surgical Robotics,'' \textit{Proc. IEEE}, vol. 110, no. 7, pp. 1012--1035, Jul. 2022.

\bibitem{ref10} W. Xia, B. Chen, S. Hu, M. Liu, M. Li, ``Torque Measurement and Safety Control Method for Joint Motors in a Harvesting Robotic Arm,'' \textit{IEEE Access}, vol. 13, pp. 213970--213985, 2025.

\bibitem{ozdamar2022} I. Ozdamar et al., ``A Shared Autonomy Reconfigurable Control Framework for Telemanipulation of Multi-Arm Systems,'' \textit{IEEE Robot. Automat. Lett.}, vol. 7, no. 4, pp. 9937--9944, Oct. 2022.

\bibitem{surgical2023} A. Noguera Cundar, R. Fotouhi, Z. Ochitwa, H. Obaid, ``Quantifying the Effects of Network Latency for a Teleoperated Compact Surgical Robot,'' 2023.

\end{thebibliography}

\section*{Anexos}

\subsection*{Repositorio de Código}
El código fuente completo del proyecto está disponible en:
\url{https://github.com/Ceciluh/Teleoperacion_con_Sensado_de_Fuerza.git}

\subsection*{Video de Demostración}
El video del sistema funcionando en tiempo real está disponible en:
\url{https://drive.google.com/file/d/1v64lXogxQqZloEJda7OgPb4155LBlYU8/view?usp=sharing }

\end{document}
%% ============================================================
%% FIN DEL DOCUMENTO
%% ============================================================
